\section{Mật mã học ứng dụng}

Trong bối cảnh các cuộc xâm nhập mạng liên tục leo thang, mã hóa không chỉ là lớp vỏ bọc mà là rào chắn sinh tồn của dữ liệu. Dự án \textit{Secret Letter} triển khai thuật toán mã hóa đối xứng AES (Advanced Encryption Standard) vận hành dưới chế độ GCM (Galois/Counter Mode), kết hợp độ dài khóa 256-bit (AES-256-GCM). Lựa chọn này không phải sự ngẫu nhiên. Nó là kết quả của việc cân bằng giữa sức mạnh tính toán mật mã và yêu cầu khắt khe về độ trễ trên môi trường trình duyệt \parencite{nist-sp800-38d-gcm}.

\subsection{Thuật toán AES và chế độ GCM}
AES đóng vai trò là một bộ mã hóa khối (block cipher), cắt nhỏ dữ liệu thành các khối 128-bit. Việc ấn định không gian khóa lên tới 256-bit là chủ ý thiết kế nhằm khắc chế tính khả thi của các cuộc tấn công vét cạn (brute-force) trên phần cứng điện toán hiện đại. 

Dù vậy, bản thân một bộ mã hóa khối thuần túy (như ECB) rất dễ bị tổn thương trước các kỹ thuật phân tích mẫu (pattern analysis). Để khắc phục, GCM được đưa vào như một mảnh ghép then chốt, mang lại cơ chế \textbf{Mã hóa xác thực kèm dữ liệu liên kết (AEAD - Authenticated Encryption with Associated Data)}. 
GCM giải quyết đồng thời hai bài toán hóc búa:
\begin{itemize}
    \item \textbf{Tính bí mật (Confidentiality):} Thay vì xử lý tuần tự, GCM biến tấu chế độ CTR (Counter) để mã hóa dữ liệu. Điều này cho phép các lõi CPU phân chia và xử lý song song (parallel processing), tối ưu hóa luồng tính toán ngay cả trên các thiết bị di động có cấu hình khiêm tốn.
    \item \textbf{Tính xác thực (Authenticity):} Dưới lăng kính của kẻ tấn công, việc dịch ngược bản mã có thể bất khả thi, nhưng chúng vẫn có thể thử chiến thuật lật bit (bit-flipping) để phá hỏng cấu trúc gói tin. GCM đập tan nỗ lực này bằng cách dùng phép nhân trên trường Galois hữu hạn (GHASH) để gắn thêm một thẻ xác thực (Authentication Tag). Bất kỳ sự xáo trộn nhỏ nào trên đường truyền cũng sẽ bị nhận diện ngay lập tức \parencite{nist-sp800-38d-gcm}.
\end{itemize}

Luồng thực thi AES-GCM đòi hỏi ba tham số đầu vào: Khóa bí mật (Secret Key), Bản rõ (Plaintext) và Vector khởi tạo (IV/Nonce). Khác với khóa mã hóa, IV không bắt buộc phải giữ bí mật nhưng tuyệt đối không được phép tái sử dụng. Dự án quy định độ dài IV là 12 byte (96-bit) và yêu cầu sinh ngẫu nhiên cho mỗi giao dịch. Sự hiện diện của IV giúp đánh lừa các kỹ thuật phân tích tần suất: cùng một đoạn thông điệp, mã hóa bằng cùng một khóa, nhưng sẽ sinh ra các khối bản mã (ciphertext) hoàn toàn riêng biệt ở mỗi lần thực thi.

\begin{figure}[H]
\centering
\begin{tikzpicture}[
    font=\sffamily\small,
    >={Stealth[round, scale=1.1]},
    box/.style={draw=blue!30, fill=blue!5, rounded corners=6pt, inner sep=8pt, align=center, minimum height=1cm, minimum width=3.5cm, font=\bfseries},
    input/.style={draw=gray!40, fill=gray!3, rounded corners=6pt, inner sep=6pt, align=center, minimum height=0.9cm, minimum width=2.8cm},
    output/.style={draw=cyan!40, fill=cyan!3, rounded corners=6pt, inner sep=6pt, align=center, minimum height=0.9cm, minimum width=2.8cm},
    arrow/.style={->, thick, gray!80, rounded corners=4pt},
    region/.style={draw=gray!30, dashed, rounded corners=12pt, inner sep=15pt, fill=gray!1}
]

% PROCESSING
\node[box] (aes) at (0,0) {Thuật toán AES-256-GCM};

% INPUTS
\node[input] (key) at (0, 2.5) {Khóa (Key)\\256-bit};
\node[input] (plain) at (-4, 2.5) {Bản rõ\\(Plaintext)};
\node[input] (iv) at (4, 2.5) {Nonce (IV)\\96-bit};

% OUTPUTS
\node[output] (cipher) at (-2.5, -2.5) {Bản mã\\(Ciphertext)};
\node[output] (tag) at (2.5, -2.5) {Thẻ xác thực\\(Auth Tag)};

% ARROWS - INPUTS to AES
\draw[arrow, out=-90, in=160, looseness=0.8] (plain.south) to (aes.160);
\draw[arrow] (key.south) -- (aes.north);
\draw[arrow, out=-90, in=20, looseness=0.8] (iv.south) to (aes.20);

% ARROWS - AES to OUTPUTS
\draw[arrow, out=200, in=90, looseness=0.8] (aes.200) to (cipher.north);
\draw[arrow, out=-20, in=90, looseness=0.8] (aes.-20) to (tag.north);

% REGIONS
\begin{scope}[on background layer]
    \coordinate (pad_in_l) at (-5.5, 2.5);
    \coordinate (pad_in_r) at (5.5, 2.5);
    \node[region, fit=(plain) (key) (iv) (pad_in_l) (pad_in_r)] (reg_in) {};
    \node[anchor=south, font=\bfseries\color{gray!80!black}] at (reg_in.north) {Dữ liệu Đầu vào (Inputs)};
    
    \coordinate (pad_out_l) at (-5.5, -2.5);
    \coordinate (pad_out_r) at (5.5, -2.5);
    \node[region, fit=(cipher) (tag) (pad_out_l) (pad_out_r)] (reg_out) {};
    \node[anchor=north, font=\bfseries\color{gray!80!black}] at (reg_out.south) {Dữ liệu Đầu ra (Outputs)};
\end{scope}

\end{tikzpicture}
\caption{Mô hình tổng quan đầu vào và đầu ra của thuật toán AES-GCM}
\label{fig:aes_gcm}
\end{figure}

\subsection{Lý do chọn AES-256-GCM}
Quyết định neo chặt kiến trúc vào AES-256-GCM xuất phát từ bốn trụ cột kỹ thuật mang tính thực tiễn cao:
\begin{enumerate}
    \item \textbf{Kháng cự lại sự ngụy tạo (Tamper Resistance):} Trong môi trường mạng phi tín nhiệm, thao túng bản mã (ciphertext malleability attack) là một chiến thuật nhắm thẳng vào tính toàn vẹn của dữ liệu. Kẻ thù có thể chặn bắt và sửa đổi vài bit ngẫu nhiên của gói tin. Khác với CBC, cơ chế Authentication Tag của GCM đóng vai trò như một người gác cổng. Hàm giải mã sẽ vứt bỏ thẳng thừng bất kỳ khối dữ liệu nào không khớp mã xác thực, hạn chế tối đa rủi ro hệ thống bị lợi dụng làm \textit{Oracle} để dò đoán bản rõ.
    \item \textbf{Tăng tốc bằng phần cứng (Hardware Acceleration):} Nhược điểm cố hữu của mã hóa Client-side là sự phụ thuộc nặng nề vào năng lực tính toán của máy khách. Dù vậy, AES-GCM được ưu ái bởi hầu hết các vi xử lý (CPU) hiện đại. Nhờ tích hợp sẵn tập lệnh chuyên biệt (chẳng hạn AES-NI) để giải quyết phép toán Galois, sự đánh đổi (trade-off) giữa bảo mật và hiệu năng được dung hòa xuất sắc, ép độ trễ mã hóa xuống mức mili-giây.
    \item \textbf{Khử rủi ro chuỗi cung ứng (Supply Chain Security):} Thay vì nhập nhằng với hàng loạt thư viện mã nguồn mở tiềm ẩn lỗ hổng (như các gói NPM), AES-GCM được trình duyệt thực thi trực tiếp qua \textit{Web Cryptography API} (\textcite{w3c-webcrypto-2017}). Việc vận hành nguyên bản (native) ở mức hệ điều hành giúp thu hẹp đáng kể bề mặt tấn công. 
    \item \textbf{Phòng ngự chiều sâu trước Điện toán lượng tử (Post-Quantum Defense):} Dưới sự đe dọa của thuật toán Grover trên máy tính lượng tử, sức mạnh của một hệ mật đối xứng sẽ bị suy giảm theo hàm căn bậc hai \parencite{nist-ir8105-pqc}. Nghĩa là một khóa 256-bit sẽ bị ép xuống mức độ khó tương đương khóa 128-bit. Mặc dù vậy, mốc 128-bit lượng tử vẫn là một bức tường thành vật lý quá sức đối với các siêu máy tính trong nhiều thập kỷ tới. Lựa chọn này giúp Secret Letter duy trì sự miễn nhiễm lâu dài trước những bước tiến của khoa học mật mã tương lai.
\end{enumerate}

Việc nhúng AES-256-GCM vào lõi hệ thống cung cấp một lời giải vẹn toàn: vừa gia tăng đáng kể độ khó cho tin tặc thông qua không gian khóa khổng lồ, vừa bảo vệ tính nguyên vẹn của luồng giao tiếp mà không vắt kiệt tài nguyên thiết bị của người dùng cuối.